\chapter{العدّ}\label{ch:b1:counting}

يبدو عدّ المجموعات \omterm{def:b1:counting:card}{المنتهية} أمرًا ابتدائيًا --- ثم يصير دقيقًا سريعًا.
يعرّف هذا الفصل \omterm{def:b1:counting:card}{عدد العناصر} تعريفًا سليمًا (عبر التقابلات، على نهج
\cref{ch:b1:logic})، ويرسي مبادئ العدّ القليلة التي يتفرع عنها
كل شيء، ثم يستخرج الأعداد الكلاسيكية: القوائم والتبديلات والأجزاء
والمعاملات الثنائية.

\section{عدد عناصر مجموعة منتهية}

\begin{definition}[المجموعة المنتهية، عدد العناصر]\label{def:b1:counting:card}
من أجل $n \in \N^*$، نكتب $\intint{1}{n} = \{1, 2, \dots,
n\}$. تكون \omterm{def:b1:logic:sets}{المجموعة} $E$ \emph{منتهية}\index{مجموعة منتهية} عندما تكون
$E = \emptyset$ أو يوجد تقابل من $\intint{1}{n}$ على $E$
من أجل $n \in \N^*$ ما؛ وهذا العدد $n$ وحيد (\cref{thm:b1:counting:welldef})
وهو \emph{عدد عناصر}\index{عدد العناصر} $E$، ويُكتب
$\abs{E}$ (مع $\abs{\emptyset} = 0$).
\end{definition}

\begin{theorem}[عدد العناصر معرَّف تعريفًا سليمًا]\label{thm:b1:counting:welldef}
إذا كان $m \neq n$ فلا يوجد تقابل من $\intint{1}{m}$
على $\intint{1}{n}$. وبدقة أكبر، إذا كان $m > n$ فلا يوجد
\omterm{def:b1:logic:inj}{تطبيق متباين} من $\intint{1}{m}$ في $\intint{1}{n}$.
\end{theorem}

\begin{proof}
نبرهن بالاستقراء على $n$ على \omterm{def:b1:logic:statement}{العبارة} التالية: \emph{لكل $m > n$
لا يوجد \omterm{def:b1:logic:inj}{تطبيق متباين} $\intint{1}{m} \to \intint{1}{n}$}. من أجل $n = 0$ تكون \omterm{def:b1:logic:sets}{مجموعة} الوصول خالية و $m \geq 1$: فلا
وجود لأيّ \omterm{def:b1:logic:map}{تطبيق} البتة. ولنفترض \omterm{def:b1:logic:statement}{العبارة} صحيحة من أجل $n$، ولنفترض أن $f \colon
\intint{1}{m} \to \intint{1}{n+1}$ \omterm{def:b1:logic:inj}{تطبيق
متباين} مع $m > n + 1$. إذا لم تُبلغ القيمة $n + 1$ كان $f$
تطبيقًا \omterm{def:b1:logic:inj}{متباينًا} في $\intint{1}{n}$، وهذا يناقض فرض
الاستقراء. وإلا فإن $f(a) = n + 1$ من أجل $a$ واحد بالضبط؛
فنبدّل بين $f(a)$ و $f(m)$ (وصوريًا: نركّب مع منقولة
القيمتين)، فيصير \omterm{def:b1:logic:inj}{التطبيق المتباين} الجديد $g$ يحقق $g(m) = n + 1$.
عندئذ يكون قصر $g$ على $\intint{1}{m-1}$ تطبيقًا
\omterm{def:b1:logic:inj}{متباينًا} في $\intint{1}{n}$ مع $m - 1 > n$ ---
وهذا تناقض من جديد.
\end{proof}

\begin{corollary}[مبدأ الأدراج]\label{cor:b1:counting:pigeonhole}
إذا كان $\abs{E} > \abs{F}$ فلا يكون أيّ \omterm{def:b1:logic:map}{تطبيق} $f \colon E \to F$
\omterm{def:b1:logic:inj}{متباينًا}: أي أن عنصرين من $E$ يشتركان في صورتهما\index{مبدأ الأدراج}.
\end{corollary}

\begin{proof}
نكتب $\abs E = m$ و $\abs F = n$ مع $m > n$، ونختار تقابلين
$u \colon \intint1m \to E$ و$v \colon F \to \intint1n$. لو كان $f$
\omterm{def:b1:logic:inj}{متباينًا} لكان $v \circ f \circ u$ تطبيقًا \omterm{def:b1:logic:inj}{متباينًا} من
$\intint1m$ في $\intint1n$ (فهو تركيب تطبيقات متباينة،
\cref{prop:b1:logic:comp})، وهذا يناقض
\cref{thm:b1:counting:welldef}.
\end{proof}

\begin{remark}[استراحة: لماذا التبديل في برهان المبرهنة؟]\label{rem:b1:counting:swapinterlude}
يحوي برهان \cref{thm:b1:counting:welldef} أولَ حركة بارعة حقًّا في
الفصل، وهي جديرة بأن تُعاد ببطء.
والعقبة هي: أننا نريد، \omterm{def:b1:logic:map}{لتطبيق} فرض الاستقراء، أن نحذف النقطة
الأخيرة $m$ من \omterm{def:b1:logic:sets}{مجموعة} التعريف \emph{و} النقطة الأخيرة
$n+1$ من \omterm{def:b1:logic:sets}{مجموعة} الوصول، لكن $f$ قد يرسل نقطة أخرى $a$ إلى $n
+ 1$، فيُفسد حذف نقطة الوصول عندئذ \omterm{def:b1:logic:map}{التطبيق} في
موضع آخر. والعلاج: أن نركّب $f$ مع منقولة \emph{القيمتين}
$f(a)$ و $f(m)$ --- وهي تقابل \omterm{def:b1:logic:sets}{لمجموعة} الوصول، فيُحفظ التباين
--- وبعدها تستقر القيمة المزعجة $n + 1$ في الموضع غير المؤذي
$m$، فيصير الحذفان نظيفين. وهذا النمط، «سوِّ أولًا ثم اقطع»،
يتكرر: فهو الذي يعيد به تراجع الاضطرابات توجيه
$\sigma^{-1}(n+1)$ في مسألة نهاية الأسبوع من هذا الفصل، وهو
كيفية ترقيع التبديلات في مسألة الزمرة المتناظرة من
\cref{ch:b1:structures}.
\end{remark}

\begin{proposition}[التطبيقات المتباينة والشاملة وعدد العناصر]\label{prop:b1:counting:injsur}
لتكن $E, F$ مجموعتين منتهيتين حيث $\abs{E} = \abs{F}$، وليكن $f \colon E
\to F$. عندئذ
\[
f \text{ متباين} \iff f \text{ شامل} \iff f \text{ تقابليّ}.
\]
\end{proposition}

\begin{proof}
نفترض $f$ \omterm{def:b1:logic:inj}{متباينًا}. عندئذ يكون $f$ تقابلًا من $E$ على $f(E)$،
ومنه $\abs{f(E)} = \abs{E} = \abs{F}$. ولو فات $f(E)$ نقطة $y_0$ من
$F$ لكان $f$ تطبيقًا \omterm{def:b1:logic:inj}{متباينًا} من $E$ في $F \setminus \{y_0\}$،
وهي \omterm{def:b1:logic:sets}{مجموعة} عدد عناصرها $\abs{F} - 1 < \abs{E}$ --- وهذا مستحيل حسب
مبدأ الأدراج. إذن $f(E) = F$: أي أن $f$ شامل، فهو إذن
تقابليّ.

ونفترض $f$ \omterm{def:b1:logic:inj}{شاملًا}. نختار من أجل كل $y \in F$ سابقةً واحدة $s(y)
\in E$؛ عندئذ $f \circ s = \mathrm{id}_F$، فيكون $s$ \omterm{def:b1:logic:inj}{متباينًا}
(\cref{prop:b1:logic:comp}). وبتطبيق الفقرة السابقة على $s$
(فعددا العناصر متساويان) يكون $s$ \omterm{def:b1:logic:inj}{تقابليًا}. ومن $f \circ s =
\mathrm{id}_F$ نحصل على $f = \mathrm{id}_F \circ s^{-1} = s^{-1}$، ومنه
يكون $f$ \omterm{def:b1:logic:inj}{تقابليًا}. وأخيرًا فإن \omterm{def:b1:logic:inj}{التطبيق التقابلي} متباين وشامل معًا
بحكم التعريف، وبهذا تُغلق دورة الاستلزامات.
\end{proof}

\begin{example}[الانتهاء شرط جوهري]\label{ex:b1:counting:finiteonly}
على \omterm{def:b1:logic:sets}{مجموعة} \emph{\omterm{def:b1:counting:card}{منتهية}} تكون \cref{prop:b1:counting:injsur} اختصارًا
قويًا: فأيّ \omterm{def:b1:logic:inj}{تطبيق متباين} من $E$ إلى نفسها هو تلقائيًا \omterm{def:b1:counting:objects}{تبديلة} في $E$ --- أي أن نصف التقابلية يأتي مجانًا. وينهار الاستلزامان معًا
على المجموعات اللانهائية: \omterm{def:b1:logic:map}{فالتطبيق} $n \mapsto n + 1$ متباين من
$\N$ إلى $\N$ لكنه يفوت $0$، \omterm{def:b1:logic:map}{والتطبيق} $\N \to \N$ الذي يرسل
$0 \mapsto 0$ و $n \mapsto n - 1$ من أجل $n \geq 1$ شامل وغير
متباين. وكلما استُدعيت هذه القضية كان فرض الانتهاء يقوم بعمل
حقيقي --- وهو موضوع تستكشفه مسألة نهاية الأسبوع في \cref{ch:b1:logic}
من الجهة المقابلة، حيث تكون المجموعات اللانهائية هي بالضبط تلك
التي تقبل تطبيقات ذاتية كهذه.
\end{example}

\begin{example}[نصف العمل، مجانًا]\label{ex:b1:counting:mod7}
تأمّل \omterm{def:b1:logic:map}{التطبيق} $f$ على $\{0, 1, \dots, 6\}$ الذي يرسل $k$ إلى
باقي قسمة $3k$ على $7$؛ وجدول قيمه هو
\[
0,\ 3,\ 6,\ 2,\ 5,\ 1,\ 4 .
\]
هل $f$ تقابل؟ يكفي التباين وحده
(\cref{prop:b1:counting:injsur}): فإذا كان للعددين $3k$ و $3k'$ الباقي نفسه قسم $7$ الفرق $3(k - k')$، وبما أن $7$ أوليّ ولا يقسم
$3$ فإنه يقسم $k - k'$ (مبرهنة إقليدس المساعدة، المستعملة هنا
على مستوى الثانوية والمبرهن عليها في \cref{ch:b1:arith})؛ ومع
$\abs{k - k'} \leq 6$ يفرض هذا $k = k'$. ويأتي الشمول
مجانًا --- دون حاجة إلى حلّ $3k \equiv c$ من أجل كل $c$، وإن
كان الجدول يؤكد ظهور كل قيمة مرة واحدة بالضبط. وهذا الاختصار
عتيد: فهو يبرهن على قابلية الضرب الترديدي للقلب
(\cref{ch:b1:arith})، وهو محرك الازدواج في مبرهنة
ويلسون، ويعود في الجبر الخطي في صورة «تشاكل ذاتي لفضاء منته
البُعد يكون \omterm{def:b1:logic:inj}{متباينًا} إذا وفقط إذا كان \omterm{def:b1:logic:inj}{شاملًا}»
(\cref{ch:b1:findim}).
\end{example}

\section{مبادئ العدّ}

\begin{proposition}[قاعدتا الجمع والجداء]\label{prop:b1:counting:rules}
لتكن $E, F$ مجموعتين منتهيتين.
\begin{enumerate}
  \item إذا كان $E \cap F = \emptyset$ فإن $\abs{E \cup F} = \abs{E} +
        \abs{F}$؛ وبعمومية أكبر، من أجل تجزئة $E$ إلى قطع
        $E_1, \dots, E_k$، لدينا $\abs{E} = \sum_i \abs{E_i}$.
  \item وفي الحالة العامة، $\abs{E \cup F} = \abs{E} + \abs{F} - \abs{E \cap
        F}$.
  \item $\abs{E \times F} = \abs{E} \times \abs{F}$.
  \item \omterm{def:b1:logic:sets}{مجموعة} كل التطبيقات من $E$ إلى $F$، وهي $F^E$، تحقق
        $\abs{F^E} = \abs{F}^{\abs{E}}$.
  \item $\abs{\mathcal{P}(E)} = 2^{\abs{E}}$.
\end{enumerate}
\end{proposition}

\begin{proof}
(1) نصل التعدادين: إذا كانت $E = \{x_1, \dots, x_m\}$ و$F =
\{y_1, \dots, y_n\}$ دون تكرار فإن $x_1, \dots, x_m, y_1,
\dots, y_n$ تعدّد $E \cup F$ دون تكرار (بحكم الانفصال).
ويمدّد الاستقراء ذلك إلى $k$ قطعة.

(2) \omterm{def:b1:logic:sets}{المجموعة} $E \cup F$ اتحاد منفصل للمجموعتين $E$ و $F \setminus E$، و
$F$ اتحاد منفصل للمجموعتين $F \cap E$ و $F \setminus E$؛ ومنه
$\abs{E \cup F} = \abs{E} + \abs{F \setminus E} = \abs{E} + \abs{F} -
\abs{E \cap F}$.

(3) \omterm{def:b1:logic:sets}{المجموعة} $E \times F$ اتحاد منفصل، على $x \in E$، للمجموعات
$\{x\} \times F$ التي \omterm{def:b1:counting:card}{عدد عناصر} كلٍّ منها $\abs{F}$؛ فطبّق (1).

(4) \omterm{def:b1:logic:map}{التطبيق} من $E = \{x_1, \dots, x_m\}$ إلى $F$ هو بالضبط اختيار
المتتالية \omterm{def:b1:counting:card}{المنتهية} ذات $m$ حدًّا $(f(x_1), \dots, f(x_m)) \in F^m$؛ وهذا
التقارن تقابل، و $\abs{F^m} = \abs{F}^m$ حسب (3)
وبالاستقراء.

(5) أجزاء $E$ تقابل تطبيقات $E \to \{0, 1\}$
(أرسل $A$ إلى دالتها المميِّزة)؛ فطبّق (4).
\end{proof}

\begin{example}[العدّ بالمتمّمة]\label{ex:b1:counting:complement}
كم عدد الرموز السرّية المكوَّنة من $4$ أرقام (الأرقام $0$--$9$،
والترتيب مهمّ، والتكرار مسموح) التي تحوي رقمًا مكررًا \emph{واحدًا
على الأقل}؟ عدّها مباشرةً يعني التلاعب بالحالات «زوج واحد بالضبط،
أو زوجان، أو ثلاثيّ، أو رباعيّ» --- وهي خمس تشكيلات متداخلة.
عُدَّ المتمّمة بدلًا من ذلك: عدد كل الرموز هو
$10^4 = 10\,000$ (بقاعدة الجداء)، وعدد الرموز ذات الأرقام الأربعة
المتمايزة هو $10 \times 9 \times 8 \times 7 = 5\,040$
(وهي الترتيبات من الرتبة $4$)، فيكون الجواب
\[
10^4 - 10 \cdot 9 \cdot 8 \cdot 7 = 10\,000 - 5\,040 = 4\,960 .
\]
أي إن قرابة نصف الرموز السرّية تكرّر رقمًا. والفكرة النافذة: كلما
صيغ عدٌّ \omterm{def:b1:logic:statement}{بعبارة} «على الأقل» أو «ليس كلها» فجرّب المتمّمة
أولًا --- إذ تضمن قاعدة الجمع أن
$\abs{A} = \abs{E} - \abs{\overline A}$، وكثيرًا ما تكون المتمّمة
تشكيلة واحدة نظيفة.
\end{example}

\begin{example}[مسارات الشبكة]\label{ex:b1:counting:paths}
عُدَّ أقصر المسارات من الزاوية $(0,0)$ إلى الزاوية
$(4, 3)$ في شبكة، بالتحرك خطوةً واحدة نحو اليمين (ي) أو خطوةً
واحدة نحو الأعلى (ف) في كل مرة. يقطع كل مسار كهذا $7$ خطوات،
منها $4$ من النوع ي و $3$ من النوع ف؛ وبالعكس، فإن أيّ كلمة طولها
$7$ من الحرفين ي و ف فيها أربعة أحرف ي تصف مسارًا واحدًا بالضبط.
فالمسارات تقابل إذن اختيارات مواضع الحروف ي:
\[
\binom{7}{4} = 35 .
\]
والفكرة النافذة هي \emph{الترميز}: فقد صار العدّ بديهيًا لحظة
تُرجم كل مسار إلى كلمة، أي إلى جزء من المواضع --- وهو مثال آخر
على الشعار القائل إن كل عدّ صحيح تقابلٌ متنكّر (\cref{met:b1:counting:which}).
\end{example}

\begin{omfigure}
\begin{tikzpicture}[scale=0.85]
  \draw[gray!60] (0,0) grid (4,3);
  \draw[->, thin] (-0.4,0) -- (4.6,0);
  \draw[->, thin] (0,-0.4) -- (0,3.6);
  \draw[omDef, very thick]
    (0,0) -- (1,0) -- (1,1) -- (3,1) -- (3,2) -- (4,2) -- (4,3);
  \fill (0,0) circle (0.07) node[below left] {$(0,0)$};
  \fill (4,3) circle (0.07) node[above right] {$(4,3)$};
  \node[omDef, below] at (0.5,0) {\small ي};
  \node[omDef, left] at (1,0.5) {\small ف};
  \node[omDef, below] at (1.5,1) {\small ي};
  \node[omDef, below] at (2.5,1) {\small ي};
  \node[omDef, left] at (3,1.5) {\small ف};
  \node[omDef, below] at (3.5,2) {\small ي};
  \node[omDef, left] at (4,2.5) {\small ف};
\end{tikzpicture}

{\small أحد المسارات القصرى، وعددها $\binom74 = 35$، من $(0,0)$ إلى
$(4,3)$: يرمّز المسار المرسوم الكلمة يفييفيف، أي اختيار
المواضع $\{1,3,4,6\}$ للحرف ي من بين الخطوات السبع.}
\end{omfigure}

\section{القوائم والتبديلات والأجزاء}

\begin{definition}[الترتيبات والتبديلات والتوفيقات]\label{def:b1:counting:objects}
لتكن $E$ \omterm{def:b1:logic:sets}{مجموعة} حيث $\abs{E} = n$ وليكن $0 \leq k \leq n$.
\begin{itemize}
  \item \emph{الترتيبة من الرتبة $k$} في $E$ هي متتالية \omterm{def:b1:counting:card}{منتهية}
        متباينة من $k$ عنصرًا من $E$ (أي اختيار مرتَّب دون تكرار)؛
  \item و\emph{التبديلة}\index{تبديلة} في $E$ هي تقابل من $E$
        إلى نفسها --- أي، على نحو مكافئ، ترتيبة من الرتبة $n$؛
  \item و\emph{التوفيقة من الرتبة $k$} هي جزء من $E$ فيه $k$
        عنصرًا (أي اختيار غير مرتَّب دون تكرار). ويُكتب عددها
        $\binom{n}{k}$، ويُقرأ «$n$ توفيق $k$»
        \index{معامل ثنائي}.
\end{itemize}
\end{definition}

\begin{theorem}[الأعداد الثلاثة]\label{thm:b1:counting:counts}
حيث $n = \abs{E}$ و $0 \leq k \leq n$:
\begin{enumerate}
  \item عدد الترتيبات من الرتبة $k$ في $E$ هو $n (n-1) \cdots
        (n-k+1) = \dfrac{n!}{(n-k)!}$؛
  \item وعدد التبديلات في $E$ هو $n!$؛
  \item و$\dbinom{n}{k} = \dfrac{n!}{k!\,(n-k)!}$.
\end{enumerate}
\end{theorem}

\begin{proof}
(1) نختار المركّبة الأولى ($n$ طريقة)، ثم الثانية ($n - 1$
اختيارًا متبقيًا)، \dots، ثم المركّبة ذات الرتبة $k$ ($n - k + 1$
اختيارًا). وصوريًا نستقرئ على $k$. من أجل $k = 1$ توجد $n$
متتالية متباينة ذات حدّ واحد. ولنفترض العدّ صحيحًا من أجل $k - 1$.
كل ترتيبة $(x_1, \dots, x_k)$ من الرتبة $k$ تُستخرج من ترتيبة
واحدة بالضبط من الرتبة $(k-1)$ --- هي بترها $(x_1, \dots, x_{k-1})$
--- بإلحاق مركّبة أخيرة خارج $\{x_1, \dots,
x_{k-1}\}$، ولها بالضبط $n - (k - 1)$ قيمة متاحة.
فتنقسم الترتيبات من الرتبة $k$ إذن، بالبتر، إلى صفوف حجمها
المشترك $n - k + 1$ مفهرسة بالترتيبات من الرتبة $(k-1)$، وتعطي
قاعدة الجمع
\[
\frac{n!}{(n-k+1)!}\;(n - k + 1) = \frac{n!}{(n-k)!} .
\]

(2) هي (1) مع $k = n$.

(3) كل جزء من $k$ عنصرًا يُرتَّب في $k!$ ترتيبة متمايزة من الرتبة
$k$، وكل ترتيبة من الرتبة $k$ تنشأ من جزء واحد بالضبط: ومنه
$\frac{n!}{(n-k)!} = \binom nk \cdot k!$.
\end{proof}

\begin{example}[الموائد المستديرة: القسمة على التناظر]\label{ex:b1:counting:roundtable}
بكم طريقة يجلس $n$ ضيفًا حول مائدة مستديرة، إذا اعتُبر جلوسان
متطابقين عندما يكون لكل ضيف الجاران نفسهما عن يمينه ويساره ---
أي بغضّ النظر عن الدوران؟ كل جلوس دائري يقابل $n$ جلوسًا خطيًا
بالضبط (اقطع الدائرة عند أيّ موضع من المواضع $n$)، فتنطوي
الترتيبات الخطية $n!$ في مجموعات من $n$:
\[
\frac{n!}{n} = (n-1)! \quad\text{جلوسًا دائريًا.}
\]
وعلى نحو مكافئ: أجلس ضيفًا مميَّزًا حيث شئت (فتقتل الحرية
الدورانية)، ثم رتّب الضيوف $n - 1$ الباقين باتجاه عقارب الساعة.
ومن أجل $n = 6$: $120$ مائدة. ويوضح الحلّان العلاجين المعياريين
للعدّ الزائد: إمّا القسمة على عدد التكرارات بالضبط، وإمّا
\emph{كسر التناظر} بتثبيت غرض واحد. ويقتضي كلاهما أن يكون حجم
زمرة التكرارات نفسه في كل تشكيلة --- وهو ما استعمله كذلك برهان
الصيغة $\binom nk = \frac{n!}{k!\,(n-k)!}$ أعلاه، مع $k!$ بدل $n$.
\end{example}

\begin{example}[إضافة قيد]\label{ex:b1:counting:constraint}
ولنواصل مع المائدة المستديرة: من بين الموائد $(n-1)!$ ذات $n \geq
3$ ضيفًا، كم مائدة تجلس ضيفين معطيين $A$ و $B$
\emph{متباعدين} (أي غير متجاورين)؟ عُدَّ المتمّمة. أمّا الموائد
التي يجلس فيها $A$ و $B$ معًا: فألصقهما في كتلة واحدة --- فيصير
لدينا $n - 1$ غرضًا حول المائدة، أي $(n-2)!$ \omterm{def:b1:logic:order}{ترتيبًا} دائريًا ---
ثم رتّب الضيفين داخل كتلتهما ($2$ طريقة): فيكون عدد الموائد
المتجاورة $2\,(n-2)!$. ومنه
\[
(n-1)! - 2\,(n-2)! = (n-2)!\,\bigl((n - 1) - 2\bigr)
= (n-3)\,(n-2)!
\]
مائدةً تُبقيهما متباعدين. وللتحقق: $n = 3$ يعطي $0$ (فحول
مثلث يلامس الجميعُ الجميعَ) و $n = 4$ يعطي $2$، ويسهل سردهما
باليد. وحيلة الإلصاق --- بمعاملة كتلة مفروضة كغرض واحد ثم عدّ
ترتيباتها الداخلية --- هي العلاج المعياري لقيود التجاور، خطيةً
كانت أو دائرية.
\end{example}

\begin{proposition}[متطابقات أساسية]\label{prop:b1:counting:identities}
من أجل $0 \leq k \leq n$:
\[
\binom{n}{k} = \binom{n}{n-k},
\qquad
\binom{n}{k} = \binom{n-1}{k-1} + \binom{n-1}{k}
\quad (1 \leq k \leq n-1),
\qquad
\sum_{k=0}^{n} \binom{n}{k} = 2^n .
\]
\end{proposition}

\begin{proof}
المتطابقة الأولى: \omterm{def:b1:logic:map}{التطبيق} $A \mapsto E \setminus A$ تقابل بين الأجزاء ذات $k$ عنصرًا والأجزاء ذات $(n-k)$ عنصرًا. وأمّا قاعدة
باسكال\index{قاعدة باسكال}: فثبّت عنصرًا $a \in E$؛ تنقسم الأجزاء
ذات $k$ عنصرًا إلى تلك التي تحتوي $a$ (فنختار العناصر $k - 1$
الأخرى: $\binom{n-1}{k-1}$) وتلك التي تتفادى $a$ ($\binom{n-1}{k}$).
وأمّا المتطابقة الثالثة: فكلا الطرفين يعدّ جميع أجزاء $E$، مقسومةً
بحسب الحجم في الطرف الأيسر
(\cref{prop:b1:counting:rules}~(1) و (5)).
\end{proof}

\begin{theorem}[مبرهنة ثنائي الحدّ]\label{thm:b1:counting:binomial}
لكل $a, b$ في حلقة تبديلية (مثل $\R$ أو $\C$) ولكل $n \in \N$:
\[
(a + b)^n = \sum_{k=0}^{n} \binom{n}{k} a^k b^{\,n-k} .
\]
\end{theorem}

\begin{proof}
نشر $(a+b)(a+b)\cdots(a+b)$ توزيعيًا يعطي حدًّا واحدًا لكل اختيار، في كل عامل، بين $a$ و $b$: فيظهر الحدّ $a^k b^{n-k}$
مرةً واحدة لكل طريقة لاختيار العوامل $k$ من بين العوامل $n$ التي تسهم بالمقدار $a$ --- أي $\binom nk$ مرة. (وبطريقة أخرى: استقرِ
على $n$ باستعمال قاعدة باسكال.)
\end{proof}

\begin{example}\label{ex:b1:counting:binomial}
تخصيصان كلاسيكيان: $a = b = 1$ يستعيد $\sum_k \binom nk
= 2^n$؛ و $a = -1$ و $b = 1$ يعطيان $\sum_{k} (-1)^k \binom nk = 0$ من أجل
$n \geq 1$: أي إن نصف أجزاء \omterm{def:b1:logic:sets}{مجموعة} غير خالية بالضبط له \omterm{def:b1:counting:card}{عدد
عناصر} زوجيّ.
\end{example}

\begin{example}[متطابقة واحدة ببرهانين]\label{ex:b1:counting:threen}
التخصيص $a = 2$ و $b = 1$ في مبرهنة ثنائي الحدّ يعطي
\[
\sum_{k=0}^{n} \binom nk\,2^k = 3^n .
\]
وإليك المتطابقة نفسها دون أيّ جبر البتة. الطرف الأيمن يعدّ
الكلمات ذات الطول $n$ على الأبجدية $\{0, 1, 2\}$
(بقاعدة الجداء). صنّف كل كلمة بحسب \omterm{def:b1:logic:sets}{مجموعة} المواضع $K$ التي
تحمل حرفًا غير معدوم: فاختيار $K$ حيث $\abs K = k$ يكلّف
$\binom nk$، ثم يحمل كل موضع من $K$ الحرف $1$ أو $2$ باستقلال:
أي $2^k$ طريقة. وتعطي قاعدة الجمع على $k$ الطرف الأيسر.
وإلى جانب متعة التوافق، لكلّ من البرهانين فضيلته: فالجبريّ
يُعمَّم على أيّ قيمة للمقدار $a$، والتوفيقيّ \emph{يشرح} الصيغة
ويتكيّف مع قيود (كمنع الحرف $2$ في الموضع الأخير مثلًا) لا
يلتقطها أيّ تعويض. والإبقاء على التقنيتين نشطتين هو المهارة
العملية التي يدرّبها هذا الفصل.
\end{example}

\begin{method}[أيّ عدّ ينطبق؟]\label{met:b1:counting:which}
قبل الحساب، أجب عن سؤالين يخصّان الاختيار: هل \emph{الترتيب}
مهمّ، وهل \emph{التكرار} مسموح؟
\begin{center}
\renewcommand{\arraystretch}{1.4}
\begin{tabular}{l|c|c}
 & الترتيب مهمّ & الترتيب غير مهمّ \\
\hline
دون تكرار & $\dfrac{n!}{(n-k)!}$ & $\dbinom{n}{k}$ \\[6pt]
بتكرار مسموح & $n^k$ & (\cref{exo:b1:counting:10}) \\
\end{tabular}
\end{center}
ثم ابحث عن تقابل أو تجزئة يردّان المسألة إلى هذه الأعداد
النموذجية؛ فكل عدّ صحيح تقابلٌ متنكّر.
\end{method}

\begin{remark}[مزالق شائعة في العدّ]\label{rem:b1:counting:pitfalls}
\begin{enumerate}
  \item \emph{جمع حالات غير منفصلة.} تقتضي قاعدة الجمع وجود
        تجزئة؛ فإذا كانت تشكيلة ما قد تحقق حالتين في آن واحد
        عُدَّت مرتين --- والعلاج هو الاحتواء والاستبعاد
        (\cref{thm:b1:counting:inclexcl}) أو تقسيم أدقّ للحالات.
  \item \emph{المرتَّب في مقابل غير المرتَّب.} اختيار «لجنة من
        اثنين» هو $\binom n2$، لا $n(n-1)$: فقرّر \emph{قبل
        الحساب} إن كان الاختيار يحمل \omterm{def:b1:logic:order}{ترتيبًا}، وإذا كان العدّ
        المرتَّب أسهل فاقسم في النهاية على عدد الترتيبات ---
        لكن بشرط أن ينشأ كل غرض غير مرتَّب من العدد \emph{نفسه}
        من الأغراض المرتَّبة.
  \item \emph{اختيارات على مراحل غير مستقلة.}
        تقتضي قاعدة الجداء أن يكون عدد الخيارات في كل مرحلة
        مستقلًا عن الاختيارات السابقة. فقولنا «اختر قائدًا ثم نائبًا
        مختلفًا عنه» سليم ($n(n-1)$)؛ أمّا «اختر لاعبين ينسجمان
        معًا» فليس جداءً على مرحلتين البتة.
  \item \emph{العدّ المزدوج بحكم الإنشاء.} بناء كل غرض مرتين ---
        مثل عدّ الأيدي التي فيها آسٌ \emph{واحد على الأقل}
        بحاصل (اختر آسًا) $\times$ (اختر $4$ أوراق أخرى) --- يعدّ
        الأيدي ذات الآسين عدًّا زائدًا. \omterm{def:b1:logic:statement}{فعبارة} «على الأقل»
        تستدعي المتمّمة في كل الأحوال تقريبًا
        (\cref{ex:b1:counting:complement}).
\end{enumerate}
\end{remark}

\begin{example}[عدٌّ على طريقة البوكر]\label{ex:b1:counting:poker}
من رزمة فيها $52$ ورقة، عدد الأيدي المكوَّنة من $5$ أوراق هو
$\binom{52}{5} = 2\,598\,960$. وأمّا الأيدي التي تحوي آسًا واحدًا بالضبط:
فاختر الآس ($4$ طرق) ثم $4$ أوراق من بين $48$ ورقة غير آس:
$4 \binom{48}{4} = 778\,320$. وتنطبق قاعدة الجداء لأن الاختيار
ينقسم إلى مراحل مستقلة.
\end{example}

\begin{method}[العدّ المزدوج]\label{met:b1:counting:doublecount}
للبرهان على متطابقة بين تعبيرين عدديين، جد \omterm{def:b1:counting:card}{مجموعة منتهية} واحدة
يعدّها الطرفان معًا --- وهي عادةً \omterm{def:b1:logic:sets}{مجموعة} من \emph{الأزواج} ---
ثم احسب عدد عناصرها بترتيبين مختلفين. والنموذج الأصلي هو
\emph{مبرهنة المصافحات} المساعدة: في حفل ما، عُدَّ الأزواج
(شخص، يد صافحها). فالجمع على الأشخاص يعطي $\sum_p d_p$ (أي عدد مصافحات كل شخص $p$)؛ والجمع على المصافحات يعطي ضعف
عدد المصافحات (إذ تضمّ كل مصافحة شخصين). ومنه فإن
$\sum_p d_p$ زوجيّ --- فعدد الأشخاص الذين صافحوا عددًا فرديًا
من الأيدي زوجيّ دائمًا، وهو استنتاج غير بديهي حصلنا عليه دون
أيّ صيغة البتة. والمحرك نفسه يشغّل
\cref{exo:b1:counting:12} وعدة أسئلة من مسألة نهاية الأسبوع
أدناه.
\end{method}

\begin{example}[الجزء المتوسط]\label{ex:b1:counting:averagesubset}
ما متوسط \omterm{def:b1:counting:card}{عدد عناصر} جزء من \omterm{def:b1:logic:sets}{مجموعة} $E$ فيها $n$ عنصرًا، إذا كانت
الأجزاء $2^n$ جميعها متساوية الاحتمال؟ عُدَّ عدًّا مزدوجًا
الأزواج $(A, a)$ حيث $a \in A$: فالجمع على الأجزاء يعطي
$\sum_A \abs A$، وهو المجموع المطلوب؛ والجمع على العناصر يعطي
$n \cdot 2^{n-1}$ (إذ يقع كلٌّ من العناصر $n$ في نصف الأجزاء
بالضبط --- فازدوج كل $A$ يحتوي $a$ مع $A
\setminus \{a\}$). ومنه
\[
\frac{1}{2^n}\sum_{A \subseteq E} \abs A
= \frac{n\,2^{n-1}}{2^n} = \frac n2 :
\]
أي إن الأجزاء نصف ممتلئة في المتوسط --- كما يتنبأ به كذلك
التناظر $A
\leftrightarrow \overline A$ (الذي يزدوج الحجمين $k$
و $n - k$). برهانان وجواب واحد، وكلاهما يتفادى الحساب المباشر
$\sum_k k\binom nk$ في \cref{exo:b1:counting:5}: فكثيرًا ما
يعوّض ازدواجٌ حسن الاختيار عن متطابقة.
\end{example}

\section{الاحتواء والاستبعاد}

\begin{theorem}[الاحتواء والاستبعاد]\label{thm:b1:counting:inclexcl}
من أجل مجموعات \omterm{def:b1:counting:card}{منتهية} $A_1, \dots, A_p$:
\[
\Bigl|\, \bigcup_{i=1}^{p} A_i \,\Bigr|
= \sum_{\emptyset \neq I \subseteq \intint{1}{p}}
(-1)^{\abs{I}+1} \Bigl|\, \bigcap_{i \in I} A_i \,\Bigr| .
\]
ومن أجل $p = 3$:
$\abs{A \cup B \cup C} = \abs A + \abs B + \abs C - \abs{A \cap B} -
\abs{A \cap C} - \abs{B \cap C} + \abs{A \cap B \cap C}$.
\end{theorem}

\begin{proof}
ثبّت عنصرًا $x$ من الاتحاد وعُدَّ إسهامه في الطرف الأيمن.
لتكن $J = \{i : x \in A_i\}$، وعدد عناصرها $m \geq
1$. يُعدّ العنصر $x$ مرة واحدة في $\abs{\bigcap_{i \in I} A_i}$
تحديدًا عندما $\emptyset \neq I \subseteq J$، بالإشارة
$(-1)^{\abs I + 1}$؛ فيكون إسهامه الكلي
\[
\sum_{k=1}^{m} \binom{m}{k} (-1)^{k+1}
= 1 - \sum_{k=0}^{m} \binom mk (-1)^k = 1 - 0 = 1
\]
حسب \cref{ex:b1:counting:binomial}. إذن يُعدّ كل عنصر من الاتحاد
مرة واحدة بالضبط.
\end{proof}

\begin{example}[عدّ الأعداد الأولية فيما بينها]\label{ex:b1:counting:coprime}
كم عددًا صحيحًا من $\intint1{120}$ يكون أوليًا مع $120 = 2^3
\times 3 \times 5$؟ يشترك عدد صحيح في عامل مع $120$ تحديدًا
عندما يقبل القسمة على $2$ أو $3$ أو $5$، فعُدَّ متمّمة
$A_2 \cup A_3 \cup A_5$، حيث تجمع $A_d$ مضاعفات
$d$. وداخل $\intint1{120}$ يكون عدد مضاعفات $d$ هو $120/d$
كلما قسم $d$ العدد $120$ --- دون حاجة إلى دوال الجزء الصحيح ---
ولدينا $A_2 \cap A_3 = A_6$ وهكذا. وبالاحتواء والاستبعاد:
\[
\abs{A_2 \cup A_3 \cup A_5}
= 60 + 40 + 24 - 20 - 12 - 8 + 4 = 88 ,
\]
ومنه فإن $120 - 88 = 32$ عددًا صحيحًا أوليّ مع $120$. ومن
المفيد إعادة تجميع الحساب في صورة جداء:
\[
120 - 88 = 120\Bigl(1 - \frac12\Bigr)\Bigl(1 -
\frac13\Bigr)\Bigl(1 - \frac15\Bigr) = 120 \cdot \frac12 \cdot
\frac23 \cdot \frac45 = 32 :
\]
فنشر الأقواس الثلاثة يعيد إنتاج الحدود المؤشَّرة الثمانية
للاحتواء والاستبعاد بالضبط، حدًّا لكل جزء من $\{2, 3,
5\}$. وهذه الصورة الجدائية تعرّف دالة أويلر المؤشِّرة، التي يظهر
دورها الحسابي مع توافقات \cref{ch:b1:arith} ويُطوَّر في
مجلد السنة 2.
\end{example}

\begin{example}[الاضطرابات]\label{ex:b1:counting:derangement}
\emph{الاضطراب} \omterm{def:b1:counting:objects}{تبديلة} بلا نقطة صامدة. لتكن $A_i$
\omterm{def:b1:logic:sets}{مجموعة} تبديلات $\intint{1}{n}$ التي تصمد عندها $i$؛
عندئذ $\abs{\bigcap_{i \in I} A_i} = (n - \abs I)!$، ويعدّ
الاحتواء والاستبعاد التبديلات ذات نقطة صامدة واحدة على الأقل؛
فيكون عدد الاضطرابات
\[
D_n = n! \sum_{k=0}^{n} \frac{(-1)^k}{k!} .
\]
وبما أن $\sum (-1)^k / k! \to \eu^{-1}$ (انظر \cref{ch:b1:series})، فإن قرابة
$37\%$ من جميع التبديلات اضطرابات، مهما تكن $n$.
\end{example}

\begin{remark}[أين يُستعمل هذا الفصل]\label{rem:b1:counting:whereused}
المعاملات الثنائية أكثر أغراض هذا الفصل إعادةَ استعمال: فهي
تحرّك مبرهنة ثنائي الحدّ في \cref{ch:b1:poly}
(نشر $(X + a)^n$)، وصيغة لايبنتز للمشتقة من الرتبة $n$ لجداء في
\cref{ch:b1:derivative}، ومعاملات نشور تايلور في \cref{ch:b1:taylor}.
وتعود التبديلات في صورة زمرة --- مع الإشارة المبنية على عدّ
الانقلابات --- في \cref{ch:b1:structures}، وتعرّف الإشارةُ بدورها
المحدداتِ في \cref{ch:b1:det}.
والاحتواء والاستبعاد ومبادئ العدّ هي العمود الفقري المنتهي
للاحتمال المتقطع، الذي يُطوَّر في مجلد السنة 2؛
أمّا أعداد الاضطراب في \cref{ex:b1:counting:derangement} فتُدرس
بعمق في مسألة نهاية الأسبوع أدناه.
\end{remark}

\section{تمارين}

\begin{exercise}[$\star$]\label{exo:b1:counting:1}
تتكوّن لوحة الترقيم من حرفين (من A إلى Z)، ثم ثلاثة أرقام، ثم
حرفين. كم لوحة ممكنة؟ وكم منها بلا حرف مكرر بين الحروف الأربعة؟
\end{exercise}

\begin{exercise}[$\star$]\label{exo:b1:counting:2}
كم قلبًا (أي إعادة ترتيب للحروف، ذا معنى أو بلا معنى)
لكلمة \textsc{برتقال}؟ وكم لكلمة \textsc{بانانا}؟
\end{exercise}

\begin{exercise}[$\star$]\label{exo:b1:counting:3}
تُختار لجنة من $4$ أشخاص من بين $7$ نساء و $5$ رجال. كم لجنة:
إجمالًا؟ وفيها $2$ من النساء بالضبط؟ وفيها رجل واحد على الأقل؟
\end{exercise}

\begin{exercise}[$\star$]\label{exo:b1:counting:4}
برهن على أنه في أيّ \omterm{def:b1:logic:sets}{مجموعة} من $13$ شخصًا يشترك اثنان في شهر
الميلاد؛ وعلى أنه من بين أيّ $n + 1$ عددًا صحيحًا مختارًا من
$\intint{1}{2n}$ يوجد عددان متتاليان. \emph{(مبدأ الأدراج في
المرتين: سمِّ الأدراج.)}
\end{exercise}

\begin{exercise}[$\star$]\label{exo:b1:counting:5}
احسب $\sum_{k=0}^{n} k \binom{n}{k}$. \emph{إرشاد: اشتقّ
$(1 + x)^n$، أو استعمل $k \binom nk = n \binom{n-1}{k-1}$ (وبرهن عليها).}
\end{exercise}

\begin{exercise}[$\star\star$]\label{exo:b1:counting:6}
كم تطبيقًا متزايدًا تمامًا يوجد من $\intint{1}{k}$ إلى $\intint{1}{n}$؟ استنتج عدد التطبيقات المتزايدة (لا
بالضرورة تمامًا). \emph{إرشاد للعدّ الثاني: \omterm{def:b1:logic:map}{التطبيق} $f$ متزايد
$\mapsto$ $g(i) = f(i) + i - 1$.}
\end{exercise}

\begin{exercise}[$\star\star$]\label{exo:b1:counting:7}
(فاندرموند) برهن، بعدّ الأجزاء ذات $k$ عنصرًا في \omterm{def:b1:logic:sets}{مجموعة} مقسومة
إلى كتلتين حجمهما $m$ و $n$، على أن:
\[
\binom{m+n}{k} = \sum_{j=0}^{k} \binom{m}{j} \binom{n}{k-j} .
\]
واستنتج $\sum_{j=0}^{n} \binom nj^2 = \binom{2n}{n}$.
\end{exercise}

\begin{exercise}[$\star\star$]\label{exo:b1:counting:8}
كم عددًا صحيحًا في $\intint{1}{1000}$ يقبل القسمة على
$2$ أو $3$ أو $5$؟ (بالاحتواء والاستبعاد؛ والمقدار $\lfloor 1000/6 \rfloor$
يعدّ مضاعفات $6$، وهكذا.)
\end{exercise}

\begin{exercise}[$\star\star$]\label{exo:b1:counting:9}
عُدَّ التطبيقات الشاملة من \omterm{def:b1:logic:sets}{مجموعة} ذات $4$ عناصر على \omterm{def:b1:logic:sets}{مجموعة} ذات
$2$ عنصرين؛ ثم على \omterm{def:b1:logic:sets}{مجموعة} ذات $3$ عناصر. \emph{إرشاد: عُدَّ
التطبيقات غير الشاملة بالاحتواء والاستبعاد على القيم الفائتة.}
\end{exercise}

\begin{exercise}[$\star\star$]\label{exo:b1:counting:10}
(النجوم والعصيّ) برهن على أن عدد الاختيارات من الرتبة $k$ من
$n$ غرضًا \emph{بتكرار} مع إهمال الترتيب --- أي، على نحو مكافئ،
عدد المتتاليات $(x_1, \dots, x_n) \in \N^n$ التي تحقق $x_1 + \dots + x_n = k$
--- هو $\binom{n + k - 1}{k}$. \emph{إرشاد: رمّز حلًّا بصفّ من
$k$ نجمة و $n - 1$ عصًا.}
\end{exercise}

\begin{exercise}[$\star\star\star$]\label{exo:b1:counting:11}
برهن على صيغة \cref{ex:b1:counting:derangement} الخاصة بالعدد $D_n$
بالتفصيل، واستنتج $n! = \sum_{k=0}^{n} \binom{n}{k} D_{n-k}$ (وبرهن
كذلك على هذه المتطابقة مباشرةً بتصنيف التبديلات بحسب \omterm{def:b1:logic:sets}{مجموعة}
نقاطها الصامدة).
\end{exercise}

\begin{exercise}[$\star\star\star$]\label{exo:b1:counting:12}
من أجل $n \in \N^*$، برهن بعدّ مزدوج للأزواج (جزء، عنصر مميَّز)
على أن:
\[
\sum_{k=1}^{n} k \binom{n}{k} = n\, 2^{n-1},
\qquad\text{ثم}\qquad
\sum_{k=1}^{n} k^2 \binom{n}{k} = n(n+1)\, 2^{n-2} .
\]
\emph{وأمّا الثانية: فعُدَّ أزواج العناصر المميَّزة، متساويةً
كانت أو مختلفة.}
\end{exercise}

\section{مسألة: الاضطرابات، أو الرسائل الخاطئة العناوين}

\begin{problem}\label{pb:b1:counting:1}
يضع سكرتير $n$ رسالة في $n$ ظرفًا معنونًا عشوائيًا:
فما احتمال ألا يتلقى \emph{أحد} رسالته الصحيحة؟
يقود هذا السؤال الكلاسيكي (مونمور، 1708) إلى أعداد الاضطراب
$D_n$ في \cref{ex:b1:counting:derangement}. وصيغة الاحتواء
والاستبعاد ليست إلا الحركة الافتتاحية: فهذه المسألة تطوّر
التراجعات التي تحسب $D_n$، وبرهانين مستقلين آخرين على الصيغة،
والمبرهنة اللافتة القائلة إن $D_n$ أقرب عدد صحيح إلى $n!/\eu$،
والتوزيع الكامل للنقاط الصامدة في \omterm{def:b1:counting:objects}{تبديلة} عشوائية، والحساب
الطريف للمتتالية $(D_n)$.\index{اضطراب}
وفي كل ما يلي، يرمز $D_n$ إلى عدد الاضطرابات (أي التبديلات
الخالية من النقاط الصامدة) في $\intint1n$، مع الاصطلاح $D_0 = 1$
(إذ إن \omterm{def:b1:counting:objects}{التبديلة} الخالية بلا نقطة صامدة).

\medskip
\noindent\textbf{الجزء 1 --- الحالات الصغيرة وإحصاء النقاط
الصامدة.}
\begin{enumerate}
  \item احسب $D_1, D_2, D_3$ مباشرةً، و $D_4$ بسرد اضطرابات
        $\{1, 2, 3, 4\}$ مجمّعةً بحسب قيمة
        $\sigma(1)$. (ينبغي أن تجد $D_4 = 9$.)
  \item من أجل $0 \leq k \leq n$، بيّن أن عدد التبديلات
        $P_k(n)$ في $\intint1n$ ذات $k$ نقطة صامدة
        \emph{بالضبط} هو $\binom nk D_{n-k}$.
  \item تحقق من الإحصاء من أجل $n = 4$: احسب $P_0(4), \dots,
        P_4(4)$ وتأكّد أن مجموعها $4! = 24$. وأيّهما أرجح من أجل
        أربع رسائل: ألا يوافق شيء، أم أن توافق رسالة واحدة بالضبط؟
  \item بعدّ مزدوج (\cref{met:b1:counting:doublecount})
        للأزواج $(\sigma, i)$ التي تحقق $\sigma(i) = i$، بيّن أن
        \[
        \sum_{\sigma} \abs{\mathrm{Fix}(\sigma)} = n! :
        \]
        أي إن \omterm{def:b1:counting:objects}{للتبديلة} العشوائية، في المتوسط، نقطة صامدة
        \emph{واحدة بالضبط}، مهما يكن $n \geq 1$.
\end{enumerate}

\medskip
\noindent\textbf{الجزء 2 --- تراجعان وبرهانان جديدان على
الصيغة.}
\begin{enumerate}[resume]
  \item برهن توفيقيًا، من أجل $n \geq 1$، على أن:
        \[
        D_{n+1} = n\,(D_n + D_{n-1}) .
        \]
        (صنّف اضطرابات $\sigma$ في $\intint1{n+1}$ بحسب
        $j = \sigma(n+1)$، ثم بحسب كون $\sigma(j) = n + 1$ أو
        لا؛ وفي حالة $\sigma(j) \neq n+1$، ابنِ تقابلًا مع
        اضطرابات $\intint1n$ بإعادة توجيه سابقة
        $n + 1$ إلى $j$.) وتحقق من التراجع عدديًا حتى
        $D_6$.
  \item بوضع $u_n = D_n - n D_{n-1}$، استنتج من السؤال 5
        أن $u_{n+1} = -u_n$، واستنتج التراجع الثاني:
        \[
        D_n = n D_{n-1} + (-1)^n \qquad (n \geq 1).
        \]
  \item انطلاقًا من السؤال 6، برهن بالاستقراء على صيغة
        \cref{ex:b1:counting:derangement}،
        \[
        D_n = n! \sum_{k=0}^{n} \frac{(-1)^k}{k!},
        \]
        --- وهو برهان مستقل كل الاستقلال عن الاحتواء والاستبعاد.
  \item (القلب الثنائي) لتكن $(a_n)$ و $(b_n)$ متتاليتين
        تحققان $a_n = \sum_{k=0}^n \binom nk b_k$ لكل
        $n$. برهن على أن
        \[
        b_n = \sum_{k=0}^{n} (-1)^{n-k} \binom nk a_k
        \qquad (n \in \N).
        \]
        (أرسِ أولًا \emph{المراجعة الثلاثية}
        $\binom nk \binom kj = \binom nj \binom{n-j}{k-j}$، ثم
        استعمل المجموع المتناوب لسطر
        \cref{ex:b1:counting:binomial}.)
  \item طبّق السؤال 8 على المتطابقة $n! = \sum_k \binom nk
        D_{n-k}$ من \cref{exo:b1:counting:11} للحصول على برهان
        \emph{ثالث} على صيغة $D_n$.
\end{enumerate}

\medskip
\noindent\textbf{الجزء 3 --- أقرب عدد صحيح إلى
$n!/\eu$.} اقبل في هذا الجزء --- فالنظرية مبنيّة في
\cref{ch:b1:series} --- أن
$\eu^{-1} = \lim_{n \to \infty} s_n$ حيث $s_n = \sum_{k=0}^{n}
\frac{(-1)^k}{k!}$، مع الحاصر التام للمتسلسلة المتناوبة
$\abs{\eu^{-1} - s_n} < \frac1{(n+1)!}$ من أجل كل $n$.
\begin{enumerate}[resume]
  \item بيّن أن $\bigl| D_n - n!/\eu \bigr| < \frac1{n+1}$ من أجل
        كل $n \in \N$.
  \item استنتج المبرهنة الرئيسة: \emph{من أجل كل $n \geq 1$،
        يكون $D_n$ أقرب عدد صحيح إلى $n!/\eu$}. ولماذا تحتاج
        الحجة إلى $n \geq 1$؟
  \item عيّن إشارة الخطأ: بيّن أن $D_n > n!/\eu$ تحديدًا عندما
        يكون $n$ زوجيًا. (حدّد موضع أول حدّ مهمَل في المتسلسلة
        المتناوبة.)
  \item احسب $D_7$ حتى $D_{10}$ بتراجع
        السؤال 5، ثم قارن $D_{10}$ بالمقدار $10!/\eu$
        ($10! = 3\,628\,800$، $\eu \approx 2.718281828$).
  \item (احتمال حفظ القبعات) لتكن $p_n = D_n/n!$ احتمال أن تكون
        \omterm{def:b1:counting:objects}{تبديلة} عشوائية منتظمة اضطرابًا. بيّن أن $\abs{p_n - \eu^{-1}} < \frac1{(n+1)!}$
        واحسب $p_6$ بخمسة أرقام عشرية. وعلّق: لماذا يكون جواب
        سؤال مونمور مستقلًا جوهريًا عن $n$ --- منذ اثنتي عشرة
        رسالة أصلًا؟
\end{enumerate}

\medskip
\noindent\textbf{الجزء 4 --- توزيع النقاط الصامدة.}
\begin{enumerate}[resume]
  \item ثبّت $k \in \N$. بيّن أن نسبة التبديلات
        في $\intint1n$ ذات $k$ نقطة صامدة بالضبط تحقق
        \[
        \frac{P_k(n)}{n!} = \frac{s_{n-k}}{k!}
        \;\xrightarrow[n \to \infty]{}\; \frac{\eu^{-1}}{k!} .
        \]
        (وهذه القيم الحدّية، ومجموعها $1$، تكوّن
        \emph{توزيع بواسون} ذا الوسيط $1$، وهو غرض محوري
        في درس الاحتمال من مجلد السنة 2.)
  \item بعدّ مزدوج للثلاثيات $(\sigma, i, j)$ حيث $i
        \neq j$ تصمد كلتاهما عند $\sigma$، بيّن أن
        $\sum_\sigma \abs{\mathrm{Fix}(\sigma)}\,
        (\abs{\mathrm{Fix}(\sigma)} - 1) = n!$ من أجل $n \geq 2$.
        وبالجمع مع السؤال 4: يكون متوسط
        $\abs{\mathrm{Fix}}^2$ مساويًا $2$، فيكون «تشتت» (تباين)
        عدد النقاط الصامدة مساويًا $1$ --- وهو مستقل عن $n$ من
        جديد، ومطابق لقانون بواسون من جديد.
  \item احسب نسبة التبديلات ذات نقطة صامدة واحدة على الأقل من
        أجل $n = 4, 5, 6$ (في صورة كسور وبأربعة أرقام عشرية)،
        وقارنها بالمقدار $1 - \eu^{-1} \approx
        0.6321$.
  \item بيّن مباشرة --- دون حاجة إلى أيّ نهايات --- أن $s_{n+2} - s_n
        = (-1)^{n+1}\bigl(\frac1{(n+1)!} - \frac1{(n+2)!}\bigr)$،
        واستنتج أن الاحتمالات $p_n = s_n$ من
        السؤال 14 تتذبذب: $p_0 > p_2 > p_4 > \dots$ و
        $p_1 < p_3 < p_5 < \dots$، فتتناقص القيم الزوجية
        وتتزايد القيم الفردية نحو النهاية المشتركة
        $\eu^{-1}$.
  \item (تبادل الهدايا السرّي) يسحب $n$ شخصًا اسمًا واحدًا من
        قبعة؛ فإذا سحب أحدهم اسمه أُعيد السحب \emph{كله} من
        جديد. باستعمال الواقعة المعيارية القائلة إن حدثًا
        احتماله $p$ يقتضي في المتوسط $1/p$ محاولة، قدّر متوسط
        عدد السحوب الكاملة اللازمة، واستنتج أن هذا الإجراء
        يكلّف نحو $\eu \approx
        2.72$ سحبة في المتوسط، وذلك باستقلال
        جوهري عن $n$.
\end{enumerate}

\medskip
\noindent\textbf{الجزء 5 --- حساب $D_n$، وتوليفة ختامية.}
\begin{enumerate}[resume]
  \item صقّل السؤال 5: بيّن أنه من أجل $j \in
        \intint2n$ مثبَّت، يكون عدد اضطرابات $\intint1n$ التي تحقق
        $\sigma(1) = j$ مساويًا $D_{n-1} + D_{n-2}$ بالضبط،
        باستقلال عن $j$. واستنتج أن $n - 1$ يقسم $D_n$
        من أجل كل $n \geq 2$.
  \item برهن على أن $D_n$ فرديّ إذا وفقط إذا كان $n$ زوجيًا.
        (اعمل بترديد $2$ في تراجع السؤال 6.)
  \item برهن على أن $D_n \equiv (-1)^n \pmod n$ من أجل $n \geq 1$،
        وتحقق من التوافق على الرقم الأخير من $D_{10}$.
  \item بيّن انطلاقًا من السؤال 6 أن $\dfrac{D_n}{D_{n-1}} = n +
        \dfrac{(-1)^n}{D_{n-1}}$ من أجل $n \geq 3$، فتكون نسبة
        عددي اضطراب متتاليين مساويةً \emph{تقريبًا بالضبط}
        $n$؛ واشرح في جملة واحدة لماذا يتوافق ذلك مع
        $D_n \approx n!/\eu$.
  \item أين استعملت هذه المسألة بالضبط: (أ) قاعدتَي الجداء
        والجمع؛ (ب) العدّ المزدوج؛ (ج) مبرهنة ثنائي
        الحدّ؛ (د) الحاصر المقبول للمتسلسلة المتناوبة؟ جملة
        واحدة لكلٍّ منها.
  \item توليفة. صار لصيغة $D_n$ الآن ثلاثة براهين
        (الاحتواء والاستبعاد، والتراجع مع الاستقراء، والقلب
        الثنائي). قارن في فقرة قصيرة ما \emph{يشرحه} كل برهان:
        أيّها أسرع حسابًا، وأيّها يتعمّم على أعداد نقاط صامدة
        أخرى، وأيّها يكشف لماذا يظهر $\eu$ في مسألة عن الأظرفة.
\end{enumerate}
\end{problem}
